\chapter{``\ldots{} Felicidade para sempre''}

\ldots{} Temos meditado, observando a nossa respiração, contemplando a inalação
e a exalação. Estamos a usar a atenção pura, consciência do corpo enquanto
caminhamos, de pé, sentados e deitados. Em vez de nos fascinarmos, estamos a
abrir a mente às condições tal como estão no presente momento.

Reparai como até num lugar bonito como este, podemos mesmo tornar a nossa vida
miserável. Quando estamos aqui, talvez queiramos estar num outro lugar qualquer;
quando estamos a andar, podemos querer estar sentados; quando estamos sentados,
talvez queiramos andar. Quando estamos a meditar, estamos a pensar no que vamos
fazer depois do retiro. E então, depois do retiro, gostaríamos de voltar\ldots{}
incorrigível, não é?

Quem sabe antes de virem para este retiro, estavam a ter problemas em casa e
pensavam: ``Mal posso esperar para ir para o retiro''. E depois aqui desejam:
``Mal posso esperar que o retiro acabe.'' Ou talvez se tranquilizem imenso aí
sentados a pensar: ``Quero estar sempre assim'', ou tentam obter aquele estado
arrebatador que obtiveram ontem, mas em vez disso irritam-se cada vez mais.

Quando obtemos estes estados agradáveis de arrebatamento, agarramo-nos a eles;
mas depois temos que arranjar algo para comer, ou fazer algo. Sentimo-nos então
mal ao perdermos o estado arrebatador. Ou talvez não tenhamos tido quaisquer
estados arrebatadores: surgem apenas muitas memórias miseráveis, raiva e
frustrações. Todas as outras pessoas estão felicíssimas; então sentimo-nos
chateados porque toda a gente parece estar a receber alguma coisa deste retiro,
excepto nós\ldots{}

É assim que começamos a observar que tudo muda. Depois temos a possibilidade de
observar como criamos problemas ou nos apegamos ao bem ou criamos todo o tipo de
complexidades em torno das condições do momento; querendo algo que não temos,
querendo manter algo que temos, querendo nos livrar de algo que temos. Este é o
problema humano do desejo, não é? Estamos sempre à procura de outra coisa.

Lembro-me de quando era criança querer um certo brinquedo. Disse à minha mãe que
se tivesse aquele brinquedo, nunca mais quereria nada. Satisfar-me-ia
completamente. E eu acreditava naquilo -- não lhe estava a dizer mentira
nenhuma; a única coisa que me impedia de realmente ser feliz naquele momento era
eu não ter o brinquedo que queria. Então a minha mãe comprou o brinquedo e
deu-mo. Consegui retirar alguma felicidade do brinquedo durante uns 5 minutos
talvez\ldots{} e então tive que começar a querer outra coisa. Portanto, ao
obter o que eu queria, senti alguma gratificação e felicidade e então o desejo
por algo mais surgia. Lembro-me disto tão vividamente porque naquela idade tão
jovem, eu realmente acreditava que se tivesse aquele brinquedo que eu queria,
seria feliz para sempre\ldots{} simplesmente para depois compreender que
``felicidade para sempre'' era uma impossibilidade\ldots{}

\photoPage{image-002-bhikkhu-meditating.jpg}{Um Bhikkhu a meditar no exterior do seu kuti}{Foto por John Johns}
